% Topic 1.3.02 — Strings and Substring Search.

\section{Строки и поиск подстроки}
\label{sec:1.3.02-strings-substring-search}

\begin{ticket}
Строки и операции над ними. Представление строк. Вычисление длины, конкатенация. Алгоритмы поиска подстроки в строке.
\end{ticket}

\subsection*{Теория}

Работаем со строками над конечным алфавитом~$\Sigma$ (например,
ASCII или кодовые точки Unicode). Основные источники~---
\cite[гл.~32]{CLRS2013} и \cite{Skiena2014}.

\subsubsection*{Строки и базовые операции}

\begin{definition}[Строка]
\label{def:string}
\emph{Строкой} (или \emph{словом}) над~$\Sigma$ называется конечная
последовательность $s = s_1 s_2 \dots s_n$ с $s_i \in \Sigma$. Её
\emph{длина} равна $|s| = n$. Единственная строка длины~$0$~---
\emph{пустая строка}, обозначается $\varepsilon$. Множество всех строк
над~$\Sigma$ обозначается $\Sigma^* = \bigcup_{n \geq 0} \Sigma^n$.
\end{definition}

\begin{definition}[Конкатенация, префикс, суффикс, подстрока]
\label{def:string-ops}
Для $s = s_1 \dots s_m$ и $t = t_1 \dots t_n$:
\begin{itemize}
  \item \emph{конкатенацией} $s t$ называется строка длины $m + n$,
        равная $s_1 \dots s_m t_1 \dots t_n$;
  \item $u$~--- \emph{префикс}~$s$ (обозначение $u \sqsubseteq s$),
        если $s = u v$ при некотором $v \in \Sigma^*$;
        \emph{собственный} префикс~--- если $u \neq s$;
  \item $u$~--- \emph{суффикс}~$s$, если $s = v u$ при некотором~$v$;
  \item $u$~--- \emph{подстрока}~$s$, если $s = v u w$ при некоторых
        $v, w$.
\end{itemize}
Множество $\Sigma^*$ с операцией конкатенации образует \emph{свободный
моноид} над~$\Sigma$ с нейтральным элементом $\varepsilon$.
\end{definition}

\subsubsection*{Представление и стоимость операций}

Два распространённых способа представить строку длины~$n$ в памяти:
\begin{itemize}
  \item \emph{С предзаписанной длиной.} Хранится явная длина~$n$ вместе
        с непрерывным массивом из $n$ символов. Вычисление $|s|$
        стоит $\Theta(1)$.
  \item \emph{С завершающим маркером (C-style).} Хранятся символы и
        затем зарезервированный маркер (например,
        \texttt{\textbackslash 0}). Вычисление $|s|$ требует
        просмотра до маркера и стоит $\Theta(n)$.
\end{itemize}

\begin{proposition}[Стоимость базовых операций]
\label{prop:string-ops-cost}
Для строк длины~$m$ и~$n$ в RAM-модели (\cref{def:complexity}):
\begin{enumerate}[label=(\alph*)]
  \item \emph{индексирование} $s_i$: $\Theta(1)$ при любом массивном
        представлении;
  \item \emph{длина} $|s|$: $\Theta(1)$ при предзаписанной длине,
        $\Theta(n)$ при завершающем маркере;
  \item \emph{конкатенация} $s t$ (с возвратом новой строки):
        $\Theta(m + n)$ по времени и памяти, поскольку каждый символ
        копируется ровно один раз;
  \item \emph{подстрока} $s_i \dots s_j$ (с возвратом новой копии):
        $\Theta(j - i + 1)$;
  \item \emph{равенство} $s = t$: $\Theta(\min\{m, n\})$ в худшем
        случае (посимвольное сравнение; различие длин~--- $O(1)$, если
        длина хранится).
\end{enumerate}
\end{proposition}
\proofof{string-ops-cost}

\begin{remark}
На практике важна стоимость многократной конкатенации. Сборка строки
длины~$n$ путём $n$ последовательных односимвольных конкатенаций над
неизменяемыми строками стоит $\Theta(n^2)$ в сумме (каждый шаг копирует
весь префикс); использование изменяемого буфера с амортизированным
удвоением снижает суммарную стоимость до $\Theta(n)$.
\end{remark}

\subsubsection*{Поиск подстроки}

\begin{definition}[Поиск подстроки]
\label{def:substring-search}
Даны \emph{текст} $T = T[1..n]$ и \emph{образец} $P = P[1..m]$ при
$m \leq n$. Требуется найти все \emph{сдвиги}
$s \in \{0, 1, \dots, n - m\}$, для которых
$T[s + 1 .. s + m] = P[1..m]$. Говорят, что $P$ \emph{встречается в}
$T$ \emph{со сдвигом~$s$}.
\end{definition}

\paragraph{Наивный алгоритм.} Перебираем все сдвиги
$s = 0, 1, \dots, n - m$; на каждом сравниваем $P$ с $T[s+1..s+m]$
посимвольно, останавливаясь на первом несовпадении.

\begin{proposition}[Сложность наивного алгоритма]
\label{prop:naive-search}
Наивный алгоритм использует $\Theta(1)$ дополнительной памяти и в
худшем случае работает за $\Theta((n - m + 1)\, m)$, то есть
$\Theta(nm)$ при $m = O(n)$. Худший случай достигается, например, на
$T = a^n$ и $P = a^{m-1} b$.
\end{proposition}
\proofof{naive-search}

\paragraph{Кнут--Моррис--Пратт (KMP).} Алгоритм KMP предварительно
вычисляет \emph{префикс-функцию}
$\pi \colon \{1, \dots, m\} \to \{0, 1, \dots, m - 1\}$ для образца, а
затем за один проход просматривает текст.

\begin{definition}[Префикс-функция]
\label{def:prefix-function}
Для строки $P[1..m]$ положим
\[
  \pi(q) \;=\; \max\{ k \in \{0, 1, \dots, q - 1\} \;:\; P[1..k] = P[q - k + 1 .. q] \}.
\]
Иначе говоря, $\pi(q)$~--- длина наибольшего \emph{собственного}
префикса $P[1..q]$, одновременно являющегося суффиксом $P[1..q]$.
\end{definition}

\begin{theorem}[Поиск KMP]
\label{thm:kmp}
Алгоритм Кнута--Морриса--Пратта находит все вхождения~$P$ в~$T$
за следующее время:
\begin{enumerate}[label=(\alph*)]
  \item $\Theta(m)$ на построение префикс-функции~$\pi$;
  \item $\Theta(n)$ на просмотр текста~$T$ при готовой~$\pi$;
  \item суммарно $\Theta(n + m)$ по времени и $\Theta(m)$ по
        дополнительной памяти.
\end{enumerate}
\end{theorem}
\proofof{kmp}

\paragraph{Рабин--Карп.} Выбираем хеш-функцию~$h$ на строках длины~$m$,
поддерживающую быстрый \emph{скользящий пересчёт}: значение
$h(T[s+2..s+m+1])$ вычисляется из $h(T[s+1..s+m])$ за $O(1)$.
Сдвигаем окно длины~$m$ по~$T$; при совпадении хеша с $h(P)$
проверяем равенство прямым сравнением.

\begin{theorem}[Рабин--Карп]
\label{thm:rabin-karp}
С полиномиальным скользящим хешем и случайным простым модулем
ожидаемое время Рабина--Карпа~--- $\Theta(n + m)$, в худшем случае~---
$\Theta(nm)$. Дополнительная память~--- $\Theta(1)$.
\end{theorem}

Мы не приводим доказательство \cref{thm:rabin-karp}; его анализ
зависит от выбора семейства хеш-функций и изложен в
\cite[\S~32.2]{CLRS2013}.

\subsection*{Доказательства}

\begin{proof}[\Cref{prop:string-ops-cost}]
\proofanchor{string-ops-cost}
\emph{(а)} Индексирование массива~--- одно обращение к памяти по
определению RAM-модели.

\emph{(б)} С предзаписанной длиной: считывание поля стоит $\Theta(1)$.
С завершающим маркером: в худшем случае просматривается вся строка;
в лучшем случае $|s| = 0$ и достаточно одного обращения.

\emph{(в)} Результат содержит $m + n$ символов, и каждый требуется
записать; меньше не получится.

\emph{(г)} Аналогично~(в) при длине $j - i + 1$.

\emph{(д)} Сравниваем символы слева направо. Первое несовпадение
завершает сравнение; в худшем случае строки совпадают вплоть до конца
более короткой, что стоит $\Theta(\min\{m, n\})$.
\end{proof}

\begin{proof}[\Cref{prop:naive-search}]
\proofanchor{naive-search}
Всего $n - m + 1$ сдвигов. На каждом сдвиге внутренний цикл
сравнения работает не более~$m$ раз, пока не найдёт несовпадения или
не подтвердит совпадение всего образца. Произведение даёт верхнюю
оценку.

Для нижней оценки берём
$T = \underbrace{aaaa\dots a}_{n}$ и
$P = \underbrace{aa\dots a}_{m - 1}\, b$. На каждом сдвиге~$s$
внутренний цикл совпадает на первых $m - 1$ символах и расходится
лишь на последнем~$b$~--- ровно~$m$ сравнений на сдвиг. Получаем
нижнюю оценку $(n - m + 1)\, m = \Theta(nm)$.
\end{proof}

\begin{proof}[\Cref{thm:kmp}]
\proofanchor{kmp}
\emph{Идея алгоритма.} Поддерживаем индекс~$q$ в~$P$, фиксирующий
длину текущего совпавшего префикса~$P$. Для каждого символа текста
$T[i]$ ($i = 1, \dots, n$):
\begin{itemize}
  \item пока $q > 0$ и $P[q + 1] \neq T[i]$, полагаем
        $q \leftarrow \pi(q)$;
  \item если $P[q + 1] = T[i]$, полагаем $q \leftarrow q + 1$;
  \item если $q = m$, фиксируем вхождение со сдвигом $i - m$ и
        полагаем $q \leftarrow \pi(q)$, чтобы продолжить поиск.
\end{itemize}
Префикс-функция~$\pi$ строится тем же правилом, применённым к самому
$P$ вместо~$T$ (со сдвигом на единицу): для $q = 2, \dots, m$ положим
$k = \pi(q - 1)$, далее, пока $k > 0$ и $P[k + 1] \neq P[q]$, полагаем
$k \leftarrow \pi(k)$; если $P[k + 1] = P[q]$, полагаем
$k \leftarrow k + 1$; в конце $\pi(q) = k$.

\emph{Корректность.} Инвариант поискового цикла: <<после обработки
$T[1..i]$ переменная~$q$ равна длине наибольшего префикса~$P$,
являющегося суффиксом $T[1..i]$>>. Инициализация ($i = 0$, $q = 0$)
тривиальна. Сохранение: по определению~$\pi$ внутренний шаг while
заменяет $q$ на следующую по убыванию длину $k < q$, для которой
$P[1..k]$ остаётся суффиксом текущей совпавшей части;
\cref{def:prefix-function} гарантирует, что таким образом
перебираются все кандидаты по убыванию. Шаг $q \leftarrow q + 1$ при
совпадении символа сохраняет инвариант.

При $q = m$ совпавший префикс имеет длину~$m$, то есть~$P$ сам
является суффиксом $T[1..i]$, и значит, $P$ встречается со
сдвигом~$i - m$.

\emph{Время работы.} На каждой итерации поискового цикла происходит
по крайней мере одно из событий: увеличение~$q$, переход $q$ к меньшему
значению через~$\pi$, или нет операций ($P[q + 1] = T[i]$ при уже
$q = 0$). Введём потенциал $\Phi = q$. Шаг увеличения стоит~$1$ и
поднимает $\Phi$ на~$1$, амортизированная стоимость~$2$. Каждый шаг
$\pi$ внутри while стоит~$1$ и уменьшает $\Phi$ хотя бы на~$1$,
амортизированная стоимость~$0$ (она оплачена предыдущим взносом). Так
как $\Phi$ начинается с~$0$ и не становится отрицательным, общая
работа по всем $n$ внешним итерациям равна $O(n)$.

То же амортизированное рассуждение для построения префикс-функции на
$m$ символах~$P$ даёт $O(m)$. Суммарное время~--- $O(n + m)$, что
совпадает с тривиальной нижней оценкой $\Omega(n + m)$ (каждый символ
требуется хотя бы один раз прочитать).
\end{proof}

\subsection*{Примеры}

\begin{example}[Прогон наивного поиска]
\label{ex:naive-search}
$T = \texttt{ababcabab}$, $P = \texttt{abab}$.
\begin{itemize}
  \item Сдвиг $s = 0$: совпадение на $T[1..4] = \texttt{abab}$.
        \emph{Совпадение.}
  \item Сдвиг $s = 1$: $P[1] = \texttt{a}$ против $T[2] = \texttt{b}$,
        несовпадение.
  \item Сдвиг $s = 2$: $P[1..2] = \texttt{ab}$ против
        $T[3..4] = \texttt{ab}$, $P[3] = \texttt{a}$ против
        $T[5] = \texttt{c}$, несовпадение.
  \item Сдвиг $s = 3$: несовпадение на первом символе.
  \item Сдвиг $s = 4$: $P[1] = \texttt{a}$ против $T[5] = \texttt{c}$,
        несовпадение.
  \item Сдвиг $s = 5$: совпадение на $T[6..9] = \texttt{abab}$.
        \emph{Совпадение.}
\end{itemize}
Итого вхождения: сдвиги $\{0, 5\}$.
\end{example}

\begin{example}[Вычисление префикс-функции]
\label{ex:prefix-function}
$P = \texttt{ababaca}$, $m = 7$:
\[
  \begin{array}{r|ccccccc}
    q       & 1 & 2 & 3 & 4 & 5 & 6 & 7 \\
    \hline
    P[q]    & a & b & a & b & a & c & a \\
    \pi(q)  & 0 & 0 & 1 & 2 & 3 & 0 & 1
  \end{array}
\]
Например, $\pi(5) = 3$: наибольший собственный префикс
$\texttt{ababa}$, являющийся также её суффиксом,~--- $\texttt{aba}$.
Далее $\pi(6) = 0$, поскольку никакой собственный префикс
$\texttt{ababac}$ не является её суффиксом. Наконец, $\pi(7) = 1$,
поскольку $\texttt{a}$~--- наибольший собственный префикс-суффикс
$\texttt{ababaca}$.
\end{example}

\begin{example}[Прогон KMP]
\label{ex:kmp-run}
С образцом $P = \texttt{ababaca}$ из предыдущего примера и текстом
$T = \texttt{abababacaba}$ алгоритм KMP проходит по каждому символу
$T$ ровно один раз и не сравнивает повторно символов, уже учтённых
префикс-функцией. Читателю предлагается убедиться, что алгоритм
сообщает о вхождении со сдвигом~$3$ ($P$ выровнен с $T[4..10]$) и
ни о каких других; общее число посимвольных сравнений~--- $11 + 7 = 18$
против $33$ у наивного алгоритма на том же входе.
\end{example}

\begin{problem}[Стоимость многократной конкатенации строк]
Допустим, строки неизменяемы: каждая конкатенация порождает свежую
копию, содержимое которой получается посимвольным копированием обоих
операндов. Пусть $f(n)$~--- сложность построения строки длины~$n$
посимвольным дописыванием в худшем случае:
$s_0 = \varepsilon$, $s_{k+1} = s_k \cdot c_{k+1}$. Покажите, что
$f(n) = \Theta(n^2)$. Затем поясните, как реализация с удвоением
буфера сводит амортизированную стоимость одного дописывания
к~$\Theta(1)$, а суммарную~--- к~$\Theta(n)$.
\end{problem}
